\item La magnificación de la imagen formada por una superficie de refracción esférica está dada por:
$$ M = -\frac{n_1 q}{n_2 p} $$
Un pisapapeles está hecho de un hemisferio de vidrio sólido con índice de refracción $1.50$. El radio de la sección transversal circular es $4.00\text{ cm}$. El hemisferio se coloca sobre su superficie plana sobre una línea dibujada de $2.50\text{ mm}$ de longitud en una hoja de papel. ¿Cuál es la longitud aparente de esta línea, vista por alguien que mira verticalmente hacia abajo a través del hemisferio?